% =============================================================================
% ISI Country Brief — Bulgaria
% External Supplier Concentration Profile
%
% ISI Paper Series — Country Brief
% International Sovereignty Institute — Research Unit
% March 2026
% =============================================================================
\documentclass[12pt,a4paper]{article}

% --- ISI Paper Series shared template ----------------------------------------
\usepackage{isi-series}

% --- PDF metadata ------------------------------------------------------------
\hypersetup{
  pdftitle={External Supplier Concentration Profile: Bulgaria — ISI EU-27 2024},
  pdfauthor={International Sovereignty Institute -- Research Unit},
  pdfcreator={International Sovereignty Institute},
  pdfsubject={ISI Paper Series -- Country Brief BG, EU-27, Vintage 2024},
  pdfkeywords={sovereignty index, EU-27, supplier concentration,
    bulgaria, country brief},
}

% ==============================================================================
\begin{document}
\hypersetup{pageanchor=false}

% ---------- Title Page with Metadata ------------------------------------------
\begin{titlepage}
\centering
\vspace*{2cm}
{\Large\textsc{ISI Paper Series --- Country Brief}}\\[1.2cm]
{\LARGE\bfseries External Supplier Concentration Profile:\\[0.3em]
Bulgaria}\\[1.5cm]
{\large International Sovereignty Index (ISI)\\[0.3em]
EU-27 --- Vintage 2024 --- Methodology v1.0}\\[1.2cm]
{\large International Sovereignty Institute --- Research Unit}\\[0.8cm]
{\large March 2026}

\vspace{0.8cm}
\noindent\rule{\textwidth}{0.4pt}
\vspace{0.4cm}

% --- Research Metadata --------------------------------------------------------
\begin{flushleft}
\noindent\textbf{Data basis:} ISI Paper~2 --- EU-27 Empirical Results 2024 (Methodology~v1.0, Vintage~2024).\\[4pt]
\noindent\textbf{Methodology:} ISI Paper~1 --- Methodological Foundations, Version~1.0.

\medskip

\noindent\textbf{JEL Codes:} F02, F15, F52, O30, Q43\\[4pt]
\noindent\textbf{Keywords:} sovereignty index; EU-27; supplier concentration; composite indicator; variance decomposition; defence dependence

\medskip

\noindent\textbf{Related publications in the ISI Paper Series:}
\begin{itemize}[nosep,leftmargin=1.5em]
\item ISI Paper~1 --- Technical Specification (v1.0) --- doi:\\
  \url{https://doi.org/10.5281/zenodo.18764227}
\item ISI Paper~2 --- EU-27 Empirical Results 2024 (v1.0) --- doi:\\
  \url{https://doi.org/10.5281/zenodo.18764170}
\end{itemize}
\end{flushleft}
\vfill
\end{titlepage}

\hypersetup{pageanchor=true}
\pagenumbering{arabic}
\pagestyle{fancy}
\fancyhead[L]{\small\sffamily\textit{ISI Country Brief}}
\fancyhead[R]{\small\sffamily\textit{Bulgaria}}
\renewcommand{\headrulewidth}{0.4pt}

% ---------- Body --------------------------------------------------------------

\section{Executive Summary}
\label{sec:summary}

Bulgaria ranks~14th in the EU-27 with a composite \ISI score of~0.340, classified as \emph{moderately concentrated}.  The dominant axes are defence~(0.861) and logistics~(0.363).  The structural profile is characterised as defence-dominant.  Under Dirichlet weight perturbation (10\,000~Monte Carlo draws), Bulgaria's mean rank is~14.15 with a standard deviation of~2.03; the 5th--95th percentile band spans ranks~11 to~18.


\section{Country Position in the EU-27}
\label{sec:position}

Bulgaria occupies rank~14 of~27 EU member states (composite score~0.340).  The EU-27 mean composite score is~0.344 with a standard deviation of~0.070.  Bulgaria's score lies below the EU-27 mean, indicating below-average external supplier concentration.  In the ranking, Bulgaria is situated between Estonia (rank~13; 0.343) and the Netherlands (rank~15; 0.335).


\section{Axis Profile}
\label{sec:axis-profile}

\Cref{tab:axis-profile} presents Bulgaria's scores on all six \ISI axes.

\begin{table}[htbp]
\centering\small
\caption{Bulgaria's \ISI axis profile.}
\label{tab:axis-profile}
\begin{tabular}{@{}lr@{}}
\toprule
\textbf{Axis} & \textbf{Score} \\
\midrule
    Finance           & 0.161 \\
    Energy            & 0.357 \\
    Technology        & 0.152 \\
    Defence           & 0.861 \\
    Critical Inputs   & 0.148 \\
    Logistics         & 0.363 \\
\bottomrule
\end{tabular}
\end{table}

The highest axis score is defence~(0.861), which lies above the EU-27 axis mean of~0.573.  The second-highest axis is logistics~(0.363), below the EU-27 axis mean of~0.518.  The lowest axis score is critical inputs~(0.148), below the EU-27 mean of~0.252.


\section{Structural Drivers}
\label{sec:drivers}

\begin{sloppypar}
Bulgaria's composite score is primarily driven by a single axis: defence~(0.861).  The gap between the highest and second-highest axis (logistics, 0.363) is~0.498, indicating a pronounced single-axis concentration.  This profile type implies that Bulgaria's rank position is sensitive to methodological choices affecting the defence axis.
\end{sloppypar}


\section{Comparative Context}
\label{sec:comparative}

Across the EU-27, the covariance-based variance decomposition shows that defence (35.1\,\%) and logistics (34.3\,\%) together account for 69.4\,\% of total cross-sectional composite variance.  Bulgaria's defence score~(0.861) lies above the EU-27 defence mean~(0.573), and its logistics score~(0.363) lies below the EU-27 logistics mean~(0.518).  Under Dirichlet weight perturbation, Bulgaria's rank standard deviation is~2.03, indicating elevated rank volatility relative to the EU-27 median.



Bulgaria's axis profile is among the most asymmetric in the EU-27.  Defence~(0.861) exceeds the second-highest axis by~0.498, and the spread between the highest and lowest axes is~0.713.  This extreme concentration on a single dimension means that Bulgaria's composite is overwhelmingly determined by one axis---the same axis that contributes most to cross-country variance.

Despite the high defence score, Bulgaria's composite~(0.340) places it near the EU-27 mean.  The explanation is arithmetical: the five remaining axes are all below their respective EU-27 means, and their combined effect offsets much of the defence elevation.  This illustrates a general property of the equal-weight composite---a country cannot rank in the top tier on the strength of one axis alone if the remaining five are subdued.

The 5th--95th~percentile band~(11--18) spans seven positions under Dirichlet perturbation.  Weight configurations that increase the defence share push Bulgaria upward into the top ten; those that dilute it allow the below-mean axes to dominate, pulling the rank toward the lower half of the distribution.  The rank standard deviation of~2.03 quantifies this sensitivity.


\section{Interpretation Boundary}
\label{sec:interpretation}

The \ISI measures concentration of external suppliers, not sovereignty in a normative or political sense.  High concentration may reflect geography, economic specialisation, or alliance structures.  A high composite score does not imply policy failure; a low score does not imply policy success.  The index provides an empirical measurement basis --- what policymakers derive from it is a separate question.


% ---------- References --------------------------------------------------------
\clearpage
\vspace*{1.5cm}

\begingroup\emergencystretch=2em\hbadness=10000
\section*{References}

\begin{enumerate}[label={[\arabic*]},leftmargin=2em,itemsep=6pt]
\item International Sovereignty Institute.
  \textit{International Sovereignty Index (ISI) technical specification, version~1.0}.
  Zenodo, 2026.\\ISI Paper Series, Paper~1.
  doi:\,\href{https://doi.org/10.5281/zenodo.18764227}{10.5281/zenodo.18764227}.

\item International Sovereignty Institute.
  \textit{International Sovereignty Index (ISI): EU-27 empirical results 2024 (v1.0)}.
  Zenodo, 2026.\\ISI Paper Series, Paper~2.
  doi:\,\href{https://doi.org/10.5281/zenodo.18764170}{10.5281/zenodo.18764170}.
\end{enumerate}
\endgroup

\end{document}
